% Feature fixture: solution and answer labels number themselves within a problem.
%
% A lone \solution / \answer / \pf lead-in, or a lone {solution} box, keeps its
% plain label ("Solution."); two or more of one kind in one worked environment
% become "Solution 1.", "Solution 2.", ... The kinds are counted apart (\answer
% and \ans together), and a box after a worked environment is numbered with the
% other boxes after it, until the next worked environment starts the count
% again. The boxes are starred so that they show in the default build the smoke
% suite runs. The lead-ins are asserted by this entry's expect= tokens in
% examples/manifest.py; the boxes, and the first pass, by
% test_solution_numbering.py.
\documentclass{didactic}

\begin{document}

\section{Lead-ins}

\begin{example}
	One solution.
	\solution LONESOLMARK
\end{example}

\begin{example}
	Two solutions.
	\solution TWOSOLAMARK
	\solution TWOSOLBMARK
\end{example}

\begin{example}
	One solution, two answers, one proof.
	\solution MIXSOLMARK
	\answer MIXANSAMARK
	\ans MIXANSBMARK
	\pf MIXPFMARK
\end{example}

\section{Boxes}

\begin{example}
	Two solution boxes and an answer box follow.
\end{example}
\begin{solution}*[1cm]
	BOXAMARK
\end{solution}
\begin{solution}*[1cm]
	BOXBMARK
\end{solution}
\begin{answer}*[1cm]
	BOXANSMARK
\end{answer}

\begin{example}
	One solution box follows.
\end{example}
\begin{solution}*[1cm]
	BOXRESTARTMARK
\end{solution}

\end{document}
